% Image credits for the photographs and AI illustrations of Book 2
% (University Chemistry, Year 1), Hindi edition. Machine-readable license
% records live in images/book2/CREDITS.md; keep this file in sync with the
% English frontmatter/image-credits-book2.tex. Authors' names stay in their
% original Latin spelling here (attribution), as in the Book 1 Hindi credits.
\chapter*{चित्र-श्रेय}

इस पुस्तक के रेखाचित्र मौलिक हैं। जहाँ छायाचित्र प्रयुक्त हैं, वे अपने-अपने मुक्त
अनुज्ञापत्रों के अंतर्गत प्रयुक्त हैं, उनके रचयिताओं के प्रति आभार सहित; हर छायाचित्र के
स्रोत की पूरी कड़ियाँ और अनुज्ञापत्रों के पते पुस्तक के रिपॉज़िटरी में
\texttt{images/book2/CREDITS.md} में सूचीबद्ध हैं।
संकट चित्रलेख रसायनों के वर्गीकरण और लेबलन की विश्व स्तर पर सुसंगत प्रणाली (GHS) के
चित्रलेख हैं, जैसे उन्हें \TeX\ Live का \texttt{ghsystem} पैकेज खींचता है।

\bigskip
छायाचित्रों और मौलिक रेखाचित्रों के साथ-साथ, कुछ चित्र
कृत्रिम बुद्धि से बनाए गए चित्रण हैं, जो पुस्तक के संपादकों द्वारा लिखे गए निर्देशों से
OpenAI के चित्र-निर्माण द्वारा बनाए गए हैं, और समावेश से पहले रासायनिक
यथार्थता के लिए जाँचे गए हैं। हर निर्मित चित्र का पूरा निर्देश
रिपॉज़िटरी में \texttt{images/book2/ai/PROMPTS.md} में अभिलिखित है।

\bigskip
\noindent छायाचित्र, अध्याय के अनुसार:
\begin{itemize}\small
  \item अध्याय~1: नील्स बोर, लगभग 1922, AB Lagrelius \& Westphal, सार्वजनिक
        अधिकार-क्षेत्र (Wikimedia Commons)।
  \item अध्याय~2: लिनस पॉलिंग, 1962, Nobel Foundation, सार्वजनिक अधिकार-क्षेत्र
        (Wikimedia Commons)।
  \item अध्याय~5: प्राकृतिक कॉपर का नमूना, Flipin, OG द्वारा, CC0 (Wikimedia
        Commons)।
  \item अध्याय~6: हैलाइट के क्रिस्टल, फ़्लुओराइट और एक अष्टफलकीय हीरा,
        James St. John द्वारा, CC BY 2.0; ग्रेफ़ाइट, Zbynek Burival द्वारा,
        CC BY-SA 4.0; हिम-क्रिस्टल, Wilson A. Bentley द्वारा (1902), सार्वजनिक
        अधिकार-क्षेत्र (सभी Wikimedia Commons)।
  \item अध्याय~8: स्वांते आरेनियस, 1909, Meisenbach Riffarth \& Co., सार्वजनिक
        अधिकार-क्षेत्र (Wikimedia Commons)।
  \item अध्याय~11: सिल्वर क्लोराइड, ब्रोमाइड और आयोडाइड के अवक्षेप,
        Rrausch1974 द्वारा, CC BY-SA 3.0 (Wikimedia Commons)।
  \item अध्याय~12: कॉपर(II) हाइड्रॉक्साइड का निलंबन, Alvy16 द्वारा, CC BY 4.0;
        कॉपर ऐम्मीन विलयन, Chemicalinterest द्वारा, सार्वजनिक अधिकार-क्षेत्र (दोनों
        Wikimedia Commons)।
  \item अध्याय~13: वाल्टर नेर्न्स्ट, 1889, Smithsonian Libraries, सार्वजनिक
        अधिकार-क्षेत्र (Wikimedia Commons)।
  \item अध्याय~16: याकोबस हेंड्रिकस वांट हॉफ़, 1911 में प्रकाशित छायाचित्र,
        सार्वजनिक अधिकार-क्षेत्र (Wikimedia Commons)।
  \item अध्याय~21: विक्टर ग्रीन्यार, सार्वजनिक अधिकार-क्षेत्र (Wikimedia Commons)।
  \item अध्याय~25: व्लादिमीर मार्कोवनिकोव, 1905 में प्रकाशित एक शोक-सूचना से लिया गया
        चित्र, सार्वजनिक अधिकार-क्षेत्र (Wikimedia Commons)।
  \item अध्याय~26: सालार दे अटाकामा के लिथियम तालाब, NASA Earth
        Observatory, सार्वजनिक अधिकार-क्षेत्र; खनिज तेल में सोडियम, Justin Urgitis द्वारा,
        CC BY-SA 2.5; ज्वाला-रंग, Hegelrast द्वारा, CC BY-SA 4.0 (सभी
        Wikimedia Commons)।
  \item अध्याय~27: पिसोलिटिक बॉक्साइट, James St. John द्वारा, CC BY 2.0; बोरैक्स,
        Aram Dulyan द्वारा, सार्वजनिक अधिकार-क्षेत्र; क्वार्ट्ज़, Stephanie Clifford द्वारा,
        CC BY 2.0; श्वेत फ़ॉस्फ़ोरस, Hi-Res Images of Chemical Elements, CC BY 3.0;
        फ़्रिट्ज़ हेबर, The Nobel Foundation, सार्वजनिक अधिकार-क्षेत्र (सभी Wikimedia
        Commons)।
  \item अध्याय~28: कावाह इजेन क्रेटर में सल्फ़र खनन, S\'emhur द्वारा,
        CC BY-SA 4.0; एम्प्यूल में क्लोरीन, W. Oelen द्वारा, CC BY-SA 3.0;
        एम्प्यूल में ब्रोमीन, Jurii द्वारा, CC BY 3.0; आयोडीन के क्रिस्टल,
        Dnn87 द्वारा, CC BY 3.0 (सभी Wikimedia Commons)।
  \item अध्याय~29: तप्त बेंच (कोफ़्लर बेंच), HBR द्वारा, CC BY-SA 3.0; ऐबे
        अपवर्तनमापी, Foreade द्वारा, CC BY-SA 4.0 (दोनों Wikimedia Commons)।
\end{itemize}
\clearpage
